\documentclass{beamer}


\title{Chapter 1 -- The Big (Bayesian) Picture}
\date{DATA 335 -- Statistical Modeling -- Winter 2026}

\begin{document}

\frame{\titlepage}

\begin{frame}
    \setlength\parskip{1em}
    \frametitle{1.1 Thinking like a Bayesian}
    Two philosophies of statistics: \textbf{Frequentist} and \textbf{Bayesian}
    
    Common goal: Learn from data about the world around us.
\end{frame}

\begin{frame}
    \setlength\parskip{1em}
    \frametitle{1.1.2 The meaning of probability}

    When flipping a fair coin, we say that \textit{the probability of flipping heads is $0.5$}.
    Which interpretation of this probability resonates with you?

    \begin{enumerate}
        \setlength\itemsep{1em}
        \item \textit{If I flip the coin over and over, I'll flip heads roughly 50\% of the time.}
        \item \textit{Heads and tails are equally plausible.}
    \end{enumerate}
\end{frame}

\begin{frame}
    \begin{enumerate}
        \setlength\parskip{1em}
        \setlength\itemsep{1em}
        \item \textit{If I flip the coin over and over, I'll flip heads roughly 50\% of the time.}
        
        
        \textbf{Frequentists} interpret probabilities as long-run frequencies of repeatable events.

        \item \textit{Heads and tails are equally plausible.}
            
        \textbf{Bayesians} interpret probabilities in terms of relative plausibilities of events.
    \end{enumerate}
\end{frame}

\begin{frame}
    \setlength\parskip{1em}
    An election is coming up and a pollster claims that candidate A has a $0.9$
    probability of winning. How do you interpret this probability?

    Which interpretation of this probability resonates with you?

    \begin{enumerate}
        \setlength\itemsep{1em}
        \item \textit{Candidate A is nine times more likely to win than to lose.}
        \item \textit{If we observe the election over and over, candidate A will win roughly 90\% of the time.}
    \end{enumerate}
\end{frame}

\begin{frame}
    \begin{enumerate}
        \setlength\parskip{1em}
        \setlength\itemsep{1em}
        \item \textit{Candidate A is nine times more likely to win than to lose.}
        
        
        \textbf{Bayesians} interpret probabilities in terms of relative plausibilities.

        \item \textit{If we observe the election over and over, candidate A will win roughly 90\% of the time.}
            
        \textbf{Frequentists} interpret probabilities as long-run frequencies of ``repeatable'' events. (Hypothetical repetitions?)
    \end{enumerate}
\end{frame}

\begin{frame}
    \setlength\parskip{1em}
    \frametitle{1.1.3 The Bayesian balancing act}

    Consider two claims:
    \begin{enumerate}
        \setlength\itemsep{1em}
        \item[(a)] Zuofu claims that he can predict the outcome of a coin.
        \item[(b)] Kavya claims that she can distinguish natural and artificial sweeteners.
    \end{enumerate}
    
    To test Zuofu's claim, you flip a fair coin 10 times and he correctly predicts all 10.
    To test Kavya's claim, you give her 10 sweetener samples and she correctly identifies each.
    In light of these experiments, what do you conclude?

    \begin{enumerate}
        \setlength\itemsep{1em}
        \item \textit{The evidence supporting Zuofu's claim is just as strong as the evidence
        supporting Kavya's claim.}
        \item \textit{You're more confident in Kavya's claim than Zuofu's claim.}
    \end{enumerate}
\end{frame}

\begin{frame}
    \begin{enumerate}
        \setlength\parskip{1em}
        \setlength\itemsep{1em}
        \item \textit{The evidence supporting Zuofu's claim is just as strong as the evidence
        supporting Kavya's claim.}
            
        In a \textbf{frequentist} analysis, ``10 out of 10'' is
        ``10 out of 10'' no matter if it's in the context of Zuofu's coins or
        Kavya's sweeteners. Strict frequentists would have be equally confidence
        in the assertions that Zuofu can predict coin flips and that Kavya can
        distinguish between natural and artificial sweetener.

        \item \textit{You're more confident in Kavya's claim than Zuofu's claim.}
        
        \textbf{Bayesian} analysis gives voice to our prior knowledge.
        Our experience on Earth suggests that Zuofu is probably overstating
        his abilities but that Kavya's claim is reasonable.
    \end{enumerate}
\end{frame}

\begin{frame}
    \setlength\parskip{1em}
    \frametitle{1.4 Asking questions}


    Suppose that during a recent doctor's visit, you tested positive for a
    rare disease. If you only get to ask the doctor one question, which
    would it be?

    \begin{enumerate}
        \setlength\itemsep{1em}
        \item \textit{What's the chance that I actually have the disease?}
        \item \textit{If I don't have the disease, what are the chances that I would've gotten this positive test result?}
    \end{enumerate}
\end{frame}

\begin{frame}
    \begin{enumerate}
        \setlength\parskip{1em}
        \setlength\itemsep{1em}
        \item \textit{What's the chance that I actually have the disease?}
            
        A \textbf{Bayesian} hypothesis test seeks to answer:
        In light of the observed data, what's the chance that the hypothesis is correct?


        \item \textit{If I don't have the disease, what are the chances that I would've gotten this positive test result?}
        
        A \textbf{frequentist} hypothesis test seeks to answer:
        If in fact the hypothesis is incorrect, what's the chance I'd have observed this, or even more extreme, data?
    \end{enumerate}
\end{frame}

\begin{frame}

    Given:
    \setlength\parskip{1em}
    \begin{itemize}
        \setlength\itemsep{1em}
        \item 2 news articles in 5 are fake.
        \item 4 in 15 fake news articles have an exclamation point in the title.
        \item 1 in 45 real news articles have an exclamation point in the title.
    \end{itemize}
    What is the probability that an article with an exclamation point
    in the title is fake?
    
\end{frame}

\begin{frame}
    \[
    P(\text{fake}\mid\text{!}) =
    \frac{P(\text{!}\mid \text{fake})P(\text{fake})}{P(\text{!})}
    \]

    \[
    P(\text{!}) = P(\text{!}\mid \text{fake})P(\text{fake}) +
    P(\text{!}\mid \text{real})P(\text{real})
    \]
\end{frame}

\begin{frame}
    \frametitle{Vampirism}
    \begin{itemize}
        \setlength\itemsep{1em}
    
        \item A new test for vampirism is on the market.
        \item A vampire tests positive for vampirism $95\%$ of the time.
        \item A mortal will test positive for vampirism $1\%$ of the time.
        \item Given that vampires make up $0.1\%$ of the population, what's the probability that someone who tests positive for vampirism is, in fact, a vampire?
    \end{itemize}
\end{frame}


\begin{frame}    
    \begin{align*}
    &P(\text{vampire}\mid+)\\
    &=\frac{P(+\mid\text{vampire})P(\text{vampire})}
    {P(+\mid\text{vampire})P(\text{vampire}) + P(+\mid\text{mortal})P(\text{mortal})}\\
    &= \frac{0.95\times 0.001}{0.95\times 0.001 + 0.01\times 0.999}\\
    &= 0.087
    \end{align*}

    Thus, an individual that tests positive for vampirism has an $8.7\%$ chance of being a vampire.
    Said differently, $91.3\%$ of positive test results are \emph{false positives}.
\end{frame}
\end{document}